\section{Introducción}

Este documento corresponde a la Especificación de Requisitos de Software (SRS) del proyecto Kit Pagos Colombia, un framework de código abierto orientado a simplificar el proceso de integración de múltiples pasarelas de pago en aplicaciones de software desarrolladas en el contexto colombiano.

El documento está organizado en las siguientes secciones principales:

\begin{itemize}
  \item \textbf{Visión del producto:} describe el propósito central del framework, su alcance y el problema que busca resolver.
  \item \textbf{Restricciones del producto:} enumera las limitaciones técnicas, regulatorias y de alcance que condicionan las decisiones de diseño e implementación.
  \item \textbf{Modelo de dominio:} presenta los conceptos centrales del dominio de pagos digitales y sus relaciones, con independencia de cualquier implementación específica.
  \item \textbf{Características de los usuarios:} identifica los perfiles de usuario que interactuarán con el framework y describe sus características relevantes.
  \item \textbf{Funciones del producto:} especifica los requisitos funcionales del sistema mediante historias de usuario y una descripción detallada de las funcionalidades esperadas por componente.
  \item \textbf{Requisitos no funcionales:} describe los atributos de calidad que el framework debe satisfacer, siguiendo la clasificación FURPS+.
\end{itemize}

El documento se apoya en el estándar IEEE 830-1998 (IEEE Recommended Practice for Software Requirements Specifications) como referencia para la estructura y el nivel de detalle de la especificación. Su audiencia principal son los miembros del equipo de desarrollo, el director del proyecto y los evaluadores académicos del trabajo de grado.
